\chapter{Central Limit Theorem}
\label{ch:11}
\section{Hermite polynomials}
\label{sec:11-1}
The centered generating function has the exact form
\[
G_N(x,v)=(1-v^2)^{N/2}\exp(x\operatorname{artanh}v).
\]
For $x=\sqrt N\xi$ and $v=t/\sqrt N$, its logarithm tends to $\xi t-t^2/2$. Consequently, for fixed $n$,
\[
N^{-n/2}k_n(\sqrt N\xi,N)\longrightarrow\operatorname{He}_n(\xi),
\]
where $\exp(\xi t-t^2/2)=\sum_n\operatorname{He}_n(\xi)t^n/n!$. When polynomial values are interpreted only on the original lattice, take admissible sequences approaching $\xi$.

\section{Discrete stochastic differential equations}
\label{sec:11-2}
The exact finite-step precursor of Hermite martingales is \eqref{eq:discrete-iterated-recursion}. In factorial normalization it reads
\[
k_n(X_{m+1},m+1)-k_n(X_m,m)
=n\,k_{n-1}(X_m,m)\epsilon_{m+1}.
\]
In particular, $X_m$, $X_m^2-m$, and $X_m^3-(3m-2)X_m$ are martingales. The update changes both the spatial variable and the order parameter; it is not the fixed-$N$ urn-chain generator of Chapter~\ref{ch:3}.

To identify the scaling, define polynomials $H_n(x,t)$ by
\[
\exp(ux-tu^2/2)=\sum_{n=0}^{\infty}H_n(x,t)\frac{u^n}{n!}.
\]
For $t>0$, $H_n(x,t)=t^{n/2}\operatorname{He}_n(x/\sqrt t)$, and coefficient differentiation gives
\[
\partial_xH_n=nH_{n-1},\qquad
\partial_tH_n+\tfrac12\partial_x^2H_n=0.
\]
For standard Brownian motion $B_t$, Ito's formula therefore yields
\[
dH_n(B_t,t)=nH_{n-1}(B_t,t)\,dB_t.
\]
This is the continuous counterpart of the displayed finite-step equation. The generating-function expansion proves fixed-degree polynomial scaling; a process-level convergence statement additionally requires a functional central limit argument and appropriate integrability control, neither of which is inferred solely from the formal similarity of the equations.

\section{Bilinear kernels and the Mehler limit}
\label{sec:11-3}
The finite foundation is the induced multinomial chain of SFA II
\cite{FeinsilverSFAII2023}, presented in
Section~\ref{sec:ency-sfa-ii}. The Gaussian and oscillator constructions
below continue that framework in a central fluctuation limit; the
cumulative-trial model is a particular stochastic realization.

The saved-success model of Section~\ref{sec:cumulative-bernoulli}
gives a finite positive-kernel derivation of the Mehler limit,
following the dynamic-occupancy viewpoint
\cite{CooperHoareRahman1977,HoareRahman1984}. Fix $0<r<1$ in its
symmetric case. Write
\[
w_N(k)=2^{-N}\binom Nk,\quad
p_{n,N}(k)=\binom Nn^{-1/2}K_n(k;N),\quad
(1+z)^{N-k}(1-z)^k=\sum_nK_n(k;N)z^n.
\]
The finite kernel relative to $w_N$ is
\begin{equation}\label{eq:ency-finite-mehler}
m_{r,N}(k,\ell)=\sum_{n=0}^Nr^np_{n,N}(k)p_{n,N}(\ell)
=\sum_{n=0}^N\frac{r^nK_n(k;N)K_n(\ell;N)}{\binom Nn}.
\end{equation}
For stationary input and output counts $K,L$,
$m_{r,N}(k,\ell)=\Pr(K=k,L=\ell)/(w_N(k)w_N(\ell))$ and
\[
\mathbb E[u^Kv^L]=
\left[\tfrac{1+r}{4}(1+uv)+\tfrac{1-r}{4}(u+v)\right]^N.
\]
Set $x_k=(N-2k)/\sqrt{2N}$ and $\delta_N=\sqrt{2/N}$.
The joint local central limit theorem gives, on compact central windows,
\[
\frac{\Pr(K=k,L=\ell)}{\delta_N^2}\longrightarrow
\frac{\exp[-(x^2-2rxy+y^2)/(1-r^2)]}{\pi\sqrt{1-r^2}},
\qquad \frac{w_N(k)}{\delta_N}\longrightarrow
w(x)=\pi^{-1/2}e^{-x^2}.
\]
Consequently
\begin{equation}\label{eq:ency-mehler-ratio}
m_{r,N}(k,\ell)\longrightarrow m_r(x,y)
=\frac1{\sqrt{1-r^2}}
\exp\left[\frac{2rxy-r^2(x^2+y^2)}{1-r^2}\right].
\end{equation}
The density is relative to $d\mu=w(x)dx$, of variance $1/2$.
The transition is $Y=rx+\sqrt{1-r^2}Z$ for $Z\sim\mu$.
For $r=e^{-t}$ its generator is $\mathcal L_{\mathrm{OU}}=\frac12D^2-xD$.
The unitary $Uf=w^{1/2}f$ conjugates $-\mathcal L_{\mathrm{OU}}$ to
$\mathcal H-\frac12I$, where $\mathcal H=\frac12(x^2-D^2)$ (the symbol $H$ is reserved for $F+G$). Thus
\[
\sqrt{w(x)}m_r(x,y)\sqrt{w(y)}
=\frac{\exp[-(1+r^2)(x^2+y^2)/(2(1-r^2))+2rxy/(1-r^2)]}
{\sqrt{\pi(1-r^2)}}
\]
is the kernel of $r^{\mathcal H-1/2}$, and multiplication by $\sqrt r$ gives
$r^{\mathcal H}$. This is Mehler's kernel in oscillator coordinates.

In the probabilists' Hermite convention of Section~\ref{sec:11-1},
the bilinear identity is
\[
\sum_{n\ge0}\frac{r^n}{n!}\operatorname{He}_n(\xi)
\operatorname{He}_n(\eta)
=\frac1{\sqrt{1-r^2}}
\exp\left[\frac{2r\xi\eta-r^2(\xi^2+\eta^2)}{2(1-r^2)}\right],
\qquad |r|<1.
\]
It converges absolutely in that disk, with the square root equal to
$1$ at $r=0$. The fixed-degree Krawtchouk--Hermite limit identifies the
spectral modes; the local-limit argument above proves convergence of
the finite kernels without an unproved interchange of an expanding
spectral sum with a limit. Neither assertion alone proves convergence
of an unrelated weighted finite-matrix spectrum.

\begin{remark}[An analytic route through the closed formula]
Section~\ref{sec:15-1} writes $m_{r,N}$ as a terminating ${}_2F_1$.
The binomial theorem converts that closed form to
$m_{r,N}(k,\ell)=[z^\ell](a+bz)^{N-k}(b+az)^k/w_N(\ell)$, with
$a=(1+r)/2$, $b=(1-r)/2$. A saddle-point evaluation of this coefficient,
with the uniform bound
$|a+be^{i\theta}|^2=|b+ae^{i\theta}|^2=1-(1-r^2)\sin^2(\theta/2)$,
gives \eqref{eq:ency-mehler-ratio} without the joint local central limit
theorem.
\end{remark}
